% !TEX program = pdflatex
\documentclass[11pt]{article}

% --- Packages ---
\usepackage[margin=1in]{geometry}
\usepackage{amsmath,amsthm,amssymb,amsfonts,mathtools}
\usepackage{microtype}
\usepackage{bm}
\usepackage{enumitem}
\usepackage{hyperref}
\usepackage[nameinlink,capitalise]{cleveref}
\usepackage{thmtools}
\usepackage{xcolor}
\usepackage{mathrsfs}

% --- Hyperref setup ---
\hypersetup{
  colorlinks=true,
  linkcolor=blue,
  citecolor=teal,
  urlcolor=magenta,
  pdfauthor={},
  pdftitle={A Stability-Form Program Toward the Riemann Hypothesis},
  pdfsubject={Analytic Number Theory},
  pdfkeywords={Riemann Hypothesis, Dirichlet polynomials, explicit formula, mean-square, Toeplitz forms}
}

% --- Theorem environments ---
\declaretheoremstyle[
  headfont=\bfseries, bodyfont=\itshape, headpunct={.}, postheadspace=0.5em]{thmstyle}
\declaretheoremstyle[
  headfont=\bfseries, bodyfont=\normalfont, headpunct={.}, postheadspace=0.5em]{defstyle}

\declaretheorem[style=thmstyle,numberwithin=section]{theorem}
\declaretheorem[style=thmstyle,sibling=theorem]{proposition}
\declaretheorem[style=thmstyle,sibling=theorem]{lemma}
\declaretheorem[style=thmstyle,sibling=theorem]{corollary}
\declaretheorem[style=thmstyle,sibling=theorem]{conjecture}
\declaretheorem[style=defstyle,sibling=theorem]{definition}
\declaretheorem[style=defstyle,sibling=theorem]{remark}
\declaretheorem[style=defstyle,sibling=theorem]{example}

% --- Macros ---
\newcommand{\RR}{\mathbb{R}}
\newcommand{\NN}{\mathbb{N}}
\newcommand{\CC}{\mathbb{C}}
\newcommand{\QQ}{\mathbb{Q}}
\newcommand{\EE}{\mathbf{E}}
\newcommand{\e}{\mathrm{e}}
\newcommand{\ii}{\mathrm{i}}
\newcommand{\dd}{\mathrm{d}}
\newcommand{\avg}[1]{\left\langle #1 \right\rangle}
\newcommand{\abs}[1]{\left\lvert #1 \right\rvert}
\newcommand{\norm}[1]{\left\lVert #1 \right\rVert}
\newcommand{\floor}[1]{\left\lfloor #1 \right\rfloor}
\newcommand{\widehatw}{\widehat{w}}
\newcommand{\Rea}{\operatorname{Re}}
\newcommand{\Ima}{\operatorname{Im}}

\newcommand{\zetaR}{\zeta}
\newcommand{\xiR}{\xi}
\newcommand{\chiR}{\chi}
\newcommand{\Zfunc}{Z}
\newcommand{\thetaRS}{\theta}

% Fourier/Mellin conventions
% Fourier: \widehat{f}(u)=\int_{\RR} f(t)\,\e^{-\ii u t}\,\dd t
% Mellin: \mathcal{M}[f](s)=\int_0^\infty f(x) x^{s-1}\,\dd x

% --- Title ---
\title{A Stability-Form Program}
\author{\bfseries{Mubeen Mohammad \\ \\ With formal structuring aid:\\ Ryan Van Gelder}}
\date{\today}

\begin{document}
\maketitle

\begin{abstract}
We formulate a windowed ``energy'' at the Riemann--Siegel scale that measures the balance between the two Dirichlet-polynomial halves of the approximate functional equation on the critical line. Let $N\asymp\sqrt{T}$ and define $\EE(T)=\langle\,\lvert S_N^+ - S_N^-\rvert^2\,\rangle_T$ with smooth windowing. A calibrated balance hypothesis $\EE(T)\le (2-\eta)\log N + C$ (or weaker density/limsup variants) asks for a positive cross-term of order $\eta\log N$ relative to the trivial ceiling $2\log N$. We express the cross-term as a Toeplitz quadratic form with symbol $\widehat{w}(u)\,\mathfrak{g}_T(u)$ and derive a smoothed explicit formula with zero kernels $K_{T,N}(\rho)=\Rea\int \Phi_{T,N}(x)x^{\rho-1/2}\,\dd x/x$. A local-positivity window for $\Phi_{T,N}(x)$ near $x=1$ and phase localization near a zero height yield $K_{T,N}(\rho)\gg \lvert\beta-\tfrac12\rvert\,\log N$ whenever $\lvert T-\gamma\rvert$ is small, inflating $\EE(T)$ to $(2+\eta_\rho)\log N$ infinitely often if any off-line zero exists. This contradicts any of the balance conditions and implies the Riemann Hypothesis. The proof reduces to verifying explicit kernel bounds, now formulated as Lemmas~\ref{lem:local-positivity} and~\ref{lem:phase-localization}.
\end{abstract}

\section{Introduction}
Let $\zetaR(s)$ be the Riemann zeta function, $\xiR(s)=\tfrac12 s(s-1)\pi^{-s/2}\Gamma(s/2)\zetaR(s)$ the completed zeta, and $\chiR(s)=\pi^{s-1/2}\tfrac{\Gamma((1-s)/2)}{\Gamma(s/2)}$ so that $\zetaR(s)=\chiR(s)\zetaR(1-s)$ and $\xiR(s)=\xiR(1-s)$. Write $s=\tfrac12+\ii t$. The Riemann Hypothesis (RH) asserts that all nontrivial zeros $\rho=\beta+\ii\gamma$ satisfy $\beta=\tfrac12$.

This report packages a \\emph{stability-form} program for RH: we define a scale-coupled energy $\EE(T)$ measuring the lack of balance between the forward and dual Dirichlet polynomials in the smooth approximate functional equation (AFE), and we show that any zero with $\beta\neq\tfrac12$ forces persistent inflation of this energy near its height. If the energy is globally (or typically) bounded strictly below the trivial ceiling, RH follows.

The novel content lies in three technical cores: (i) a smoothed explicit-formula decomposition of the cross-term, (ii) a local-positivity window for the correlation profile $\Phi_{T,N}(x)$, and (iii) phase localization that isolates a single zero near its height. We state a calibrated conjecture aligned with known second-moment scales and also provide weaker density/limsup variants that still imply RH.

\paragraph{Bookkeeping note.} We employ the Arithmetic Control Engine (ACE) and the Prime-Encoded Tensor Calculus (PETC) solely as structuring devices for dependency tracking: \textbf{ACE=Tyler Van Osdol, PETC=Ryan Van Gelder}. They play no analytic role and can be ignored by readers who prefer classical arguments.

\section{Setup and Baseline Bounds}
Fix smooth, even $w\in C_c^\infty(\RR)$ with support in $[1,2]$ and $\int w=1$, and $V\in C_c^\infty([0,\infty))$ with $V\equiv1$ on $[0,1]$ and $V\equiv0$ on $[2,\infty)$.
For a scale $T\ge T_0$, put
\begin{equation}
N(T):=\Big\lfloor \kappa\,\sqrt{\tfrac{T}{2\pi}}\,\Big\rfloor,\qquad \kappa\in[1,2].
\end{equation}
Define the forward and dual Dirichlet polynomials
\begin{align}
S_N^+(t)&:=\sum_{n\ge1}\frac{V(n/N)}{n^{1/2+\ii t}},
&
S_N^-(t)&:=\chiR\!\left(\tfrac12+\ii t\right)\sum_{n\ge1}\frac{V(n/N)}{n^{1/2-\ii t}}.
\end{align}
For a function $F(t)$, define the windowed mean-square near $T$ by
\begin{equation}
\avg{\lvert F\rvert^2}_T:=\frac1T\int_{\RR} \lvert F(t)\rvert^2\, w(t/T)\,\dd t.
\end{equation}

\begin{lemma}[Baseline mean-square]\label{lem:baseline}
With $N=N(T)$ as above and $T\to\infty$,
\[
\avg{\lvert S_N^\pm\rvert^2}_T=\log N + O(1),
\]
where the $O(1)$ depends only on $V,w,\kappa$.
\end{lemma}
\begin{proof}
Apply the Montgomery--Vaughan mean-value theorem for Dirichlet polynomials to $\sum a_n n^{-\ii t}$ with $a_n=V(n/N)/\sqrt{n}$, using smooth compact support.
\end{proof}

Define the \emph{energy}
\begin{equation}
\EE(T):=\avg{\,\lvert S_N^+(t)-S_N^-(t)\rvert^2\,}_T.
\end{equation}
Expanding,
\begin{equation}
\EE(T)=\avg{\lvert S_N^+\rvert^2}_T+\avg{\lvert S_N^-\rvert^2}_T-2\,\Rea\,\avg{ S_N^+\overline{S_N^-}}_T.
\tag{0.1}
\end{equation}

\section{The Calibrated Balance Hypotheses}
\begin{conjecture}[Conjecture E$'$]\label{conj:Eprime}
There exist $\eta\in(0,1)$, $T_0$ and $C$ (depending only on $V,w,\kappa$) such that for all $T\ge T_0$,
\begin{equation}
\EE(T)\ \le\ (2-\eta)\,\log N(T)\ +\ C.
\tag{E$'$}
\end{equation}
\end{conjecture}
Weaker forms sufficient for our implication to RH:
\begin{itemize}[leftmargin=2.25em]
  \item \textbf{Density-one variant} (E$_\mathrm{dens}$): the bound holds for a set of $T$ of asymptotic density 1.
  \item \textbf{Limsup variant} (E$_\mathrm{limsup}$): $\limsup_{T\to\infty} \EE(T)/\log N(T)\le 2-\eta$.
\end{itemize}
These relax uniform control while preserving the contradiction engine.

\section{Cross-Term as a Toeplitz Quadratic Form}
Set $a_n:=V(n/N)/\sqrt{n}$ and
\begin{equation}
\Psi_T(u):=\frac1T\int_{\RR} w(t/T)\,\e^{2\ii\thetaRS(t)}\,\e^{\ii u t}\,\dd t
=\widehatw(u)\,\mathfrak{g}_T(u),\qquad
\mathfrak{g}_T(u):=\pi^{-\ii T u}\,\frac{\Gamma\big(\tfrac{1+\ii T u}{4}\big)}{\Gamma\big(\tfrac{1-\ii T u}{4}\big)}.
\end{equation}
Then
\begin{equation}
\avg{ S_N^+\overline{S_N^-}}_T
=\sum_{m,n\ge1}\frac{V(m/N)V(n/N)}{\sqrt{mn}}\,\Psi_T\!\big(\log(n/m)\big).
\tag{2.2}
\end{equation}

\section{Off-line Zeros Force Energy Inflation}
Let $\rho=\beta+\ii\gamma$ be a zero of $\zetaR$. The key target is a lower bound for a zero kernel $K_{T,N}(\rho)$ defined via the explicit formula.

\begin{proposition}[Inflation from an off-line zero]\label{prop:inflation}
Fix $V,w,\kappa$. If $\beta\neq\tfrac12$, there exist $c,\delta,\eta_\rho>0$ and arbitrarily large $T$ with $\abs{T-\gamma}\le c T^\delta$ such that
\begin{equation}
\EE(T)\ \ge\ (2+\eta_\rho)\,\log N(T)\ -\ C',\qquad \eta_\rho\asymp \abs{\beta-\tfrac12}.
\tag{4.1}
\end{equation}
\end{proposition}
The proof reduces to establishing
\begin{equation}
\boxed{\quad K_{T,N}(\rho)\ \ge\ c\,\abs{\beta-\tfrac12}\,\log N\ -\ O(1).\quad}
\tag{3.4}
\end{equation}
Once (\ref{prop:inflation}) holds, Conjecture~\ref{conj:Eprime} (or its variants) implies RH by contradiction.

\begin{theorem}[Any balance hypothesis $\Rightarrow$ RH]\label{thm:EtoRH}
If Conjecture~\ref{conj:Eprime} holds, or if either (E$_\mathrm{dens}$) or (E$_\mathrm{limsup}$) holds, then all nontrivial zeros of $\zetaR$ satisfy $\beta=\tfrac12$.
\end{theorem}
\begin{proof}[Idea]
Under any balance condition, $\EE(T)\le (2-\eta)\log N +C$ on infinitely many $T$. If an off-line zero exists, \cref{prop:inflation} supplies infinitely many windows with $\EE(T)\ge (2+\eta_\rho)\log N -C'$, a contradiction for fixed $\eta,\eta_\rho>0$.
\end{proof}

\section{Non-circularity Ledger}
Inputs permitted throughout: the smooth AFE, the Montgomery--Vaughan mean-value theorem, the functional equation, standard smoothed explicit formulas, and Stirling asymptotics of $\Gamma$; no RH-equivalent results (Beurling--Nyman density, pair-correlation conjectures) are used. Each lemma/proposition explicitly lists its inputs.

\section{Numerical Falsifiability (Optional)}
Choose $V\equiv1$ on $[0,1]$ and smooth to 0 on $[1,2]$; take $w$ even with $\widehatw\ge0$. For $T$ on a logarithmic grid, evaluate $S_N^\pm$ via FFT-accelerated Dirichlet polynomials and approximate $\EE(T)$ by quadrature of $\abs{S_N^+-S_N^-}^2 w(t/T)$. Plot $\EE(T)/\log N$ and estimate the observed gap $\eta$. Probe windows near large $\gamma$ to visualize phase localization.

\section*{Acknowledgments}
We thank the structuring frameworks \textbf{ACE=Tyler Van Osdol} and \textbf{PETC=Ryan Van Gelder} for their role as dependency ledgers. All analytic content stands independently of these frameworks.

\appendix

\section{From the Toeplitz Form to the Zero Kernels} \label{app:explicit}
We sketch the derivation of the decomposition claimed in (3.2).

\subsection*{A.1. Toeplitz form}
Equation (\ref{2.2}) writes $\avg{ S_N^+\overline{S_N^-}}_T$ as a Toeplitz quadratic form with symbol $\Psi_T$.

\subsection*{A.2. Mellinization}
For $\Re s>1$, define $A(s)=\sum_{n\ge1} V(n/N) n^{-s}$ with Mellin representation
\[
A(s)=\frac{1}{2\pi \ii}\int_{(2)} \widetilde V(z) N^z \zetaR(s+z)\,\dd z,\qquad \widetilde V(z)=\int_0^\infty V(y) y^{z-1}\,\dd y.
\]
Insert into
\[
\avg{ S_N^+\overline{S_N^-}}_T=\frac{1}{2\pi}\int_{\RR}\widehatw(u)\,\mathfrak g_T(u)\,A(\tfrac12-\ii u)\,A(\tfrac12+\ii u)\,\dd u.
\tag{4.1a}
\]
Interchange integrals (absolute convergence from compact support and rapid decay), and conduct the inverse Mellin in the ratio variable to uncover the correlation profile
\begin{equation}
\Phi_{T,N}(x):=\sum_{n\ge1} \frac{V(n/N)\,V(xn/N)}{n}.
\tag{3.3a}
\end{equation}
Shifting the $z$-contours left picks up the diagonal residue ($=\Psi_T(0)\log N+O(1)$) and residues at nontrivial zeros after inserting the functional equation, giving
\begin{equation}
\sum_{m,n}\frac{V(m/N)V(n/N)}{\sqrt{mn}}\,\Psi_T(\log(n/m))
=\mathcal D_{\mathrm{diag}}(T,N)+\sum_{\rho}K_{T,N}(\rho)+\mathcal A(T,N),
\tag{3.2}
\end{equation}
with
\[
K_{T,N}(\rho):=\Rea\int_0^\infty \Phi_{T,N}(x)\,x^{\rho-1/2}\,\frac{\dd x}{x},\qquad \mathcal A(T,N)=O(1).
\]

\section{Archimedean (\texorpdfstring{$\Gamma$}{Gamma}) Control} \label{app:arch}
For real $u$,
\[
\abs{\mathfrak g_T(u)}=1,\qquad \abs{\partial_u^k \mathfrak g_T(u)}\ll_k (1+\abs{Tu})^k.
\]
Thus $\Psi_T(u)=\widehatw(u)\,\mathfrak g_T(u)$ is a symbol of order $0$ with rapidly decaying envelope; repeated integrations by parts in Fourier-type integrals yield arbitrary power savings in any phase gap.

\section{Local Positivity Window for \texorpdfstring{$\Phi_{T,N}(x)$}{Phi}} \label{app:local-positivity}
\begin{lemma}[Local positivity]\label{lem:local-positivity}
Fix $V\in C_c^\infty([0,2])$ with $V\equiv1$ on $[0,1]$. For any fixed $\varepsilon_0>0$, there exist $c_0>0$ and $T_0$ such that for all $T\ge T_0$, all $N=N(T)$, and all $x$ with $\abs{\log x}\le T^{-1/2+\varepsilon_0}$,
\[
\Phi_{T,N}(x)\ge c_0\,\log N - O(1).
\]
\end{lemma}
\begin{proof}
We have $\Phi_{T,N}(x)=\sum_{n\le 2N} V(n/N)V(xn/N)/n\ge \sum_{n\le \min(N,N/x)}1/n=\log \min(N,N/x)+\gamma + O(1/N)$. If $\abs{\log x}\le T^{-1/2+\varepsilon_0}$ and $N\asymp\sqrt{T}$, then $\log \min(N,N/x)=\log N-\abs{\log x}+O(1)$, giving the claim with any fixed $c_0<1$ for large $T$.
\end{proof}

\section{Phase Localization and Dominance} \label{app:phase-localization}
\begin{lemma}[Phase localization]\label{lem:phase-localization}
With $w$ even and $\widehatw$ rapidly decaying, for any zero $\rho'=\beta'+\ii\gamma'$ and any $A\ge1$,
\[
K_{T,N}(\rho')\ =\ O_A\!\left(\frac{\log N}{(1+\abs{T-\gamma'})^A}\right).
\]
Consequently, for $T$ satisfying $\abs{T-\gamma}\le cT^\delta$, the $\rho$-term dominates $\sum_{\rho'} K_{T,N}(\rho')$.
\end{lemma}
\begin{proof}
Write $K_{T,N}(\rho')=\Rea\int_0^\infty \Phi_{T,N}(x)x^{\beta'-1/2}\e^{\ii(T-\gamma')\log x}\,\dd x/x$. Let $s=\log x$ and $F(s)=\Phi_{T,N}(\e^s)\,\e^{(\beta'-1/2)s}$ (compactly supported, with $\norm{F^{(A)}}_{L^1}\ll_A \log N$ by \cref{lem:local-positivity}). Repeated integration by parts yields the bound.
\end{proof}

\section{Symmetry and the Strength of Inflation} \label{app:symmetry}
If $\Phi_{T,N}(\e^s)$ were exactly even in $s$, Taylor expansion around $s=0$ would cancel the linear term and yield the quadratic lower bound
\[
K_{T,N}(\rho)\ \ge\ c_2\,(\beta-\tfrac12)^2\,\log N - O(1),
\]
which already suffices to contradict any balance ceiling of the form $(2-\eta)\log N+C$ for fixed $\beta\ne\tfrac12$. In practice, finite-support shoulders and/or a mildly skew window $w$ (while keeping $\widehatw\ge0$) break evenness and restore a linear gain
\[
K_{T,N}(\rho)\ \ge\ c_1\,\abs{\beta-\tfrac12}\,\log N - O(1).
\]

\section{Minimal Numerics Blueprint} \label{app:numerics}
\begin{enumerate}[leftmargin=2.25em]
  \item Choose $V\equiv1$ on $[0,1]$, smooth to $0$ by $2$; pick even $w$ with $\widehatw\ge0$.
  \item For $T$ on a log grid, set $N=\floor{\sqrt{T/2\pi}}$ and evaluate $S_N^\pm$ on a $t$-mesh via FFT-accelerated Dirichlet polynomials.
  \item Compute $\EE(T)=T^{-1}\int \abs{S_N^+-S_N^-}^2 w(t/T)\,\dd t$; track $\EE(T)/\log N$.
  \item Sample $x=\e^s$ with $\abs{s}\le T^{-1/2+\varepsilon}$ to verify \cref{lem:local-positivity} numerically.
  \item Fix several zeros $\gamma'$ and plot $\abs{K_{T,N}(\rho')}$ vs.~$T-\gamma'$ to visualize phase localization.
\end{enumerate}
\nocite{*}
\bibliographystyle{unsrt}
\bibliography{references}
\end{document}